% AA-35 | Audyt historyczny V: pelne uzasadnienie podnoszenia Frobeniusa.
\chapter[Redukcja modulo p i cykle Frobeniusa]{Redukcja modulo liczb pierwszych i cykle Frobeniusa}\label{aa:ch:frobenius-reduction}

\section{Motywacja i dobra redukcja}
Niech $f\in\mathbb Z[X]$ będzie monicznym wielomianem rozdzielczym.
W ciele rozkładu nad $\mathbb Q$ trudno zwykle wypisać automorfizmy
jego pierwiastków. Znacznie łatwiej rozłożyć $f$ modulo niewielką
liczbę pierwszą $p$. Pokażemy, że stopnie czynników takiego rozkładu
wyznaczają typ cykli \emph{pewnego} elementu grupy Galois nad
$\mathbb Q$. Ważne jest zastrzeżenie o dobrej redukcji:
pierwiastki nie mogą się ze sobą zlać po redukcji modulo $p$.

\begin{definicja}\label{aa:def:good-reduction}
Niech $f\in\mathbb Z[X]$ będzie monicznym wielomianem stopnia
$n\ge1$, o niezerowym wyróżniku $\operatorname{disc}(f)$.
Liczbę pierwszą $p$ nazywamy \emph{dobrą dla $f$}, jeśli
$p\nmid\operatorname{disc}(f)$.
Przez $\bar f\in\mathbb F_p[X]$ oznaczamy wielomian powstały
przez redukcję każdego współczynnika $f$ modulo $p$.
Ponieważ $f$ jest moniczny, $\deg\bar f=\deg f$.
\end{definicja}

\begin{lemat}\label{aa:lem:good-squarefree}
Niech $f\in\mathbb Z[X]$ będzie moniczny stopnia $n\ge1$
i niech jego wyróżnik będzie niezerowy. Dla dowolnej liczby
pierwszej $p$ następujące warunki są równoważne:
\begin{enumerate}
\item $p$ jest liczbą pierwszą dobrą dla $f$;
\item $\gcd(\bar f,\bar f')=1$ w $\mathbb F_p[X]$;
\item $\bar f$ nie ma pierwiastka wielokrotnego w domknięciu
algebraicznym $\mathbb F_p$.
\end{enumerate}
\end{lemat}
\begin{proof}
Najpierw uzasadnimy zgodność wyróżnika z redukcją współczynników.
Wielomian
$D(X_1,\ldots,X_n)=\prod_{i<j}(X_i-X_j)^2$
jest symetryczny i ma współczynniki całkowite. Podstawowe twierdzenie
o wielomianach symetrycznych
(Twierdzenie~\ref{aa:tw:fundamental-symmetric}), zastosowane nad
$\mathbb Z$, daje $D=P(e_1,\ldots,e_n)$ z pewnym
$P\in\mathbb Z[T_1,\ldots,T_n]$.
Dla każdego monicznego wielomianu
$g=X^n+c_1X^{n-1}+\cdots+c_n$ nad dowolnym ciałem wzory Viète'a
wyrażają jego wyróżnik jako
$P(-c_1,c_2,\ldots,(-1)^nc_n)$.
Ponieważ $P$ ma współczynniki całkowite, otrzymujemy
\[
 \overline{\operatorname{disc}(f)}
 =\operatorname{disc}(\bar f)\quad\text{w }\mathbb F_p.
\]
Z definicji dobrej redukcji i z kryterium wyróżnika
(Twierdzenie~\ref{aa:tw:disc-separable}) wynika zatem
równoważność warunków (1) i (3).

Pozostaje porównać (2) i (3). Niech $\alpha$ będzie pierwiastkiem
$\bar f$ w domknięciu algebraicznym $\overline{\mathbb F}_p$.
Piszemy $\bar f=(X-\alpha)^m h$, gdzie $m\ge1$ i
$h(\alpha)\ne0$. Z reguły różniczkowania iloczynu wynika,
że gdy $m=1$, to $(\bar f)'(\alpha)=h(\alpha)\ne0$,
a gdy $m\ge2$, to $(\bar f)'(\alpha)=0$.
Zatem $\bar f$ i $(\bar f)'$ mają wspólny pierwiastek
w domknięciu algebraicznym dokładnie wtedy, gdy $\bar f$
ma pierwiastek wielokrotny. Nad ciałem jest to równoważne
temu, że ich największy wspólny dzielnik ma dodatni stopień.
Rozumowanie obejmuje także przypadek $(\bar f)'=0$:
ponieważ $\deg\bar f\ge1$, wówczas
$\gcd(\bar f,0)=\bar f$ nie jest jednostką, a każdy
pierwiastek $\bar f$ jest wielokrotny.

Nie korzystaliśmy z nieuprawnionego utożsamiania
macierzy Sylvestera o różnych rozmiarach: po redukcji
modulo $p$ stopień pochodnej może bowiem się zmniejszyć.
\end{proof}

\section{Cykle Frobeniusa nad ciałami skończonymi}
\begin{twierdzenie}[Cykle Frobeniusa]\label{aa:tw:frob-cycles}
Niech $q=p^r$, gdzie $p$ jest liczbą pierwszą, a $r\ge1$.
Niech $h\in\mathbb F_q[X]$ będzie monicznym, rozdzielczym
wielomianem stopnia $n\ge1$, o rozkładzie
$h=h_1\cdots h_s$ na parami różne moniczne czynniki
nierozkładalne, $\deg h_i=d_i$.
Niech $E$ będzie ciałem rozkładu $h$ nad $\mathbb F_q$.
Wtedy automorfizm $\varphi:E\to E$, $x\mapsto x^q$,
działa na $n$ pierwiastkach $h$ jako iloczyn rozłącznych
cykli długości $d_1,\ldots,d_s$.
\end{twierdzenie}
\begin{proof}
Wybierzmy $i\in\{1,\ldots,s\}$ i pierwiastek $\alpha_i$
wielomianu $h_i$. Ponieważ $h_i$ jest nierozkładalny,
$[\mathbb F_q(\alpha_i):\mathbb F_q]=d_i$.
Z twierdzenia o ciałach skończonych
(Twierdzenie~\ref{aa:tw:finite-frobenius-roots})
$\mathbb F_q(\alpha_i)$ ma $q^{d_i}$ elementów,
a wszystkie jego elementy spełniają $x^{q^{d_i}}=x$.
Stąd $\varphi^{d_i}(\alpha_i)=\alpha_i$.

Pokażemy, że żaden dodatni $k<d_i$ nie spełnia
$\varphi^k(\alpha_i)=\alpha_i$. Gdyby tak było,
$\alpha_i$ należałby do zbioru pierwiastków
$X^{q^k}-X$, czyli do ciała $\mathbb F_{q^k}$.
W ustalonym domknięciu algebraicznym
$\mathbb F_q(\alpha_i)$ byłoby wtedy podciałem
jednoznacznie określonego ciała $\mathbb F_{q^k}$.
Ze wzoru na stopnie wieży otrzymalibyśmy
$d_i\mid k$, co jest niemożliwe dla $0<k<d_i$.
Zatem orbita $\alpha_i$ pod iteracjami $\varphi$
ma dokładnie $d_i$ elementów.

Wszystkie te elementy są pierwiastkami $h_i$:
automorfizm Frobeniusa ustala współczynniki z $\mathbb F_q$.
Ponieważ $h_i$ jest rozdzielczy i ma dokładnie $d_i$
pierwiastków, orbita jest całym zbiorem jego pierwiastków.
Różne czynniki $h_i$ mają rozłączne zbiory pierwiastków,
więc odpowiadające im cykle również są rozłączne.
\end{proof}

\section{Podnoszenie Frobeniusa do grupy Galois nad $\mathbb Q$}
W następnym dowodzie wykorzystamy pierścień generowany
przez \emph{pierwiastki rozpatrywanego wielomianu}.
Dzięki temu nie musimy zakładać bez dowodu struktury całego
pierścienia liczb całkowitych ciała rozkładu.

\begin{lemat}[Skończony pierścień generowany przez pierwiastki]
\label{aa:lem:frob-finite-order}
Niech $f\in\mathbb Z[X]$ będzie monicznym wielomianem
stopnia $n\ge1$, a $\alpha_1,\ldots,\alpha_n$ jego
pierwiastkami w pewnym ciele rozkładu $L/\mathbb Q$,
liczonymi z krotnościami. Wtedy pierścień
$B=\mathbb Z[\alpha_1,\ldots,\alpha_n]$ jest skończenie
generowany jako $\mathbb Z$-moduł, każdy jego element jest
całkowity nad $\mathbb Z$, a $B\cap\mathbb Q=\mathbb Z$.
Jeżeli $f$ jest rozdzielczy i $G=\operatorname{Gal}(L/\mathbb Q)$,
to $B$ jest stabilny względem $G$ i $B^G=\mathbb Z$.
\end{lemat}
\begin{proof}
Każdy $\alpha_i$ spełnia moniczne równanie
$\alpha_i^n+c_{n-1}\alpha_i^{n-1}+\cdots+c_0=0$
z całkowitymi współczynnikami. Pozwala ono zastępować
potęgi $\alpha_i^k$ dla $k\ge n$ kombinacjami całkowitymi
niższych potęg. Stosując to kolejno do wszystkich zmiennych,
otrzymujemy, że $B$ jest generowany nad $\mathbb Z$ przez
najwyżej $n^n$ monomów
$\alpha_1^{e_1}\cdots\alpha_n^{e_n}$, $0\le e_i<n$.

Niech $b\in B$, a $z_1,\ldots,z_N$ będzie skończoną rodziną
generatorów $B$ jako $\mathbb Z$-modułu.
Ponieważ $b z_i\in B$, istnieje macierz
$C=(c_{ij})\in M_N(\mathbb Z)$ taka, że
$b z_i=\sum_j c_{ij}z_j$. Dla wektora kolumnowego
$z=(z_1,\ldots,z_N)^T$ zapisujemy to jako $bz=Cz$.
Niech $P(T)=\det(TI_N-C)\in\mathbb Z[T]$.
Z twierdzenia Cayleya–Hamiltona mamy $P(C)=0$,
a z równości $b^kz=C^kz$ wynika $P(b)z=0$.
Ponieważ $z_i$ generują $B$ i $1\in B$, dostajemy $P(b)=0$.
Wielomian $P$ jest moniczny, więc $b$ jest całkowity
nad $\mathbb Z$.

Jeśli $b=u/v\in B\cap\mathbb Q$, gdzie $u,v\in\mathbb Z$,
$\gcd(u,v)=1$ i $v>0$, to z monicznego równania
$b^d+c_{d-1}b^{d-1}+\cdots+c_0=0$ po pomnożeniu
przez $v^d$ wynika $v\mid u^d$. Względna pierwszość
$u,v$ wymusza $v=1$. Stąd $B\cap\mathbb Q=\mathbb Z$.

Jeśli $f$ jest rozdzielczy, każdy $\sigma\in G$ permutuje
$\alpha_1,\ldots,\alpha_n$, a zatem zachowuje $B$.
Elementy stałe działania $G$ w $L$ tworzą $\mathbb Q$,
więc $B^G=B\cap\mathbb Q=\mathbb Z$.
\end{proof}

\begin{lemat}[Podnoszenie Frobeniusa]
\label{aa:lem:frob-lift}
Niech $L/\mathbb Q$ będzie ciałem rozkładu monicznego,
rozdzielczego wielomianu $f\in\mathbb Z[X]$, niech
$G=\operatorname{Gal}(L/\mathbb Q)$ i
$B=\mathbb Z[\alpha_1,\ldots,\alpha_n]$ jak wyżej.
Dla dowolnej liczby pierwszej $p$ i dowolnego ideału
maksymalnego $\mathfrak p$ pierścienia $B$, zawierającego $p$,
istnieje $\sigma\in G$ spełniające
$\sigma(\mathfrak p)=\mathfrak p$ i
\[
 \overline{\sigma(b)}=\bar b^{\,p}
 \quad\text{dla każdego }b\in B,
\]
gdzie kreska oznacza klasę modulo $\mathfrak p$.
\end{lemat}
\begin{proof}
\emph{Krok 1: skończoność ciał reszt.}
Z Lematu~\ref{aa:lem:frob-finite-order} pierścień $B$
jest skończonym $\mathbb Z$-modułem. Ponieważ $1/p$ nie jest
całkowity nad $\mathbb Z$, nie może należeć do $B$;
zatem $pB\ne B$. Iloraz $B/pB$ jest niezerowym skończonym
pierścieniem, gdyż jest generowany jako grupa addytywna
przez skończenie wiele elementów rzędu dzielącego $p$.
Każdy jego ideał maksymalny ma iloraz będący skończonym
ciałem. W szczególności
$\kappa=B/\mathfrak p$ jest skończonym rozszerzeniem $\mathbb F_p$.
W pierścieniu $B$ nad $p$ leży tylko skończenie wiele
ideałów maksymalnych,
bo odpowiadają one ideałom maksymalnym skończonego pierścienia
$B/pB$.

\emph{Krok 2: zastosowanie chińskiego twierdzenia o resztach.}
Dla różnych ideałów maksymalnych $\mathfrak p_i$, $\mathfrak p_j$
mamy $\mathfrak p_i+\mathfrak p_j=B$.
Wybierzmy dla każdego $j\ne i$ element $t_{ij}\in\mathfrak p_j$
taki, że $t_{ij}\equiv1\pmod{\mathfrak p_i}$.
Iloczyn $e_i=\prod_{j\ne i}t_{ij}$, równy $1$, gdy nie ma
innych ideałów, jest równy $1$ modulo $\mathfrak p_i$
i $0$ modulo każdego innego ideału.
Kombinacje liniowe $e_i$ z wybranymi reprezentantami klas
pokazują, że można zadawać niezależnie reszty modulo wszystkie
ideały maksymalne nad $p$.

\emph{Krok 3: przechodniość działania $G$ na ideałach nad $p$.}
Gdyby pewien ideał maksymalny $\mathfrak q\supset pB$
nie należał do orbity $\mathfrak p$, to
$\mathfrak p\ne\sigma^{-1}(\mathfrak q)$ dla każdego $\sigma\in G$.
Z Kroku 2 można wybrać $b\in B$ takie, że
$b\equiv0\pmod{\mathfrak p}$ oraz
$b\equiv1\pmod{\sigma^{-1}(\mathfrak q)}$
dla wszystkich $\sigma\in G$.
Norma $N=\prod_{\sigma\in G}\sigma(b)$ należy do
$B^G=\mathbb Z$. Ponieważ $b\in\mathfrak p$, także
$N\in\mathfrak p\cap\mathbb Z=p\mathbb Z\subset\mathfrak q$.
Jednocześnie dla każdego $\sigma\in G$ zachodzi
$\sigma(b)\notin\mathfrak q$, gdyż
$b\notin\sigma^{-1}(\mathfrak q)$.
Ideał maksymalny jest pierwszy, więc ich iloczyn
$N\notin\mathfrak q$, sprzeczność. Zatem działanie jest przechodnie.

\emph{Krok 4: podniesienie automorfizmu Frobeniusa.}
Niech $D=\{\sigma\in G:\sigma(\mathfrak p)=\mathfrak p\}$.
Wybierzmy element $c\in\kappa^\times$ generujący
multiplikatywną grupę $\kappa^\times$; wówczas
$\kappa=\mathbb F_p(c)$. Dzięki Krokowi 2 istnieje $b\in B$
o klasie $c$ modulo $\mathfrak p$ i klasie $0$ modulo
każdego innego ideału nad $p$.
Rozpatrzmy
$P(T)=\prod_{\sigma\in G}(T-\sigma(b))\in\mathbb Z[T]$.
Współczynniki tego wielomianu należą do $B$ i są stałe
względem $G$, a więc należą do $B^G=\mathbb Z$.
W $\kappa$ element $c$ jest pierwiastkiem redukcji $P$,
bo dla czynnika z $\sigma=1$ mamy $c-\bar b=0$.
Zatem $c^p$ również jest pierwiastkiem redukcji $P$,
ponieważ jej współczynniki należą do $\mathbb F_p$.
Wobec tego $c^p=\overline{\sigma(b)}$ dla pewnego $\sigma\in G$.
Ponieważ $c^p\ne0$, mamy $b\notin\sigma^{-1}(\mathfrak p)$.
Wszystkie ideały $\sigma^{-1}(\mathfrak p)$ są maksymalne i leżą nad $p$;
z naszego wyboru $b$ wynika więc
$\sigma^{-1}(\mathfrak p)=\mathfrak p$, czyli $\sigma\in D$.
Ponieważ $\sigma(\mathfrak p)=\mathfrak p$, działanie $\sigma$
na $B$ indukuje dobrze określony automorfizm ciała
$\kappa=B/\mathfrak p$; jego odwrotność indukuje
$\sigma^{-1}$, a więc jest on bijektywny. Automorfizm ten
ustala $\mathbb F_p$ i wysyła generator $c$ na $c^p$.
Każdy element $\kappa=\mathbb F_p(c)$ jest funkcją wymierną
$c$ o współczynnikach w $\mathbb F_p$, więc automorfizm ten
zgadza się na całym $\kappa$ z Frobeniusem
$x\mapsto x^p$. To daje żądane $\sigma$ dla wszystkich $b\in B$.
\end{proof}

\begin{twierdzenie}[Redukcja i grupa Galois]\label{aa:tw:dedekind-cycle}
Niech $f\in\mathbb Z[X]$ będzie moniczny, rozdzielczy,
stopnia $n\ge1$, a $L$ jego ciałem rozkładu nad $\mathbb Q$.
Niech $p$ będzie dowolną dobrą liczbą pierwszą dla $f$.
Jeżeli w $\mathbb F_p[X]$ wielomian $\bar f$ rozkłada się
na moniczne czynniki nierozkładalne stopni
$d_1,\ldots,d_s$, to $\operatorname{Gal}(L/\mathbb Q)$
zawiera element, którego działanie na $n$ różnych
pierwiastkach $f$ ma rozkład na rozłączne cykle
długości $d_1,\ldots,d_s$.
\end{twierdzenie}
\begin{proof}
Weźmy pierścień $B$ z Lematu~\ref{aa:lem:frob-finite-order}
i ideał maksymalny $\mathfrak p$ zawierający $pB$;
istnieje on, gdyż $B/pB$ jest niezerowym skończonym pierścieniem.
Z Lematu~\ref{aa:lem:frob-lift} istnieje $\sigma\in G$
stabilizujące $\mathfrak p$ i działające na
$\kappa=B/\mathfrak p$ jako Frobenius $x\mapsto x^p$.

Wszystkie $\alpha_i$ należą do $B$.
Gdyby dwie ich redukcje $\bar\alpha_i,\bar\alpha_j$
przy $i\ne j$ były równe, to $\alpha_i-\alpha_j\in\mathfrak p$.
W iloczynie
$\operatorname{disc}(f)=\prod_{i<j}(\alpha_i-\alpha_j)^2$
znalazłby się czynnik z $\mathfrak p$, więc
$\operatorname{disc}(f)\in\mathfrak p\cap\mathbb Z=p\mathbb Z$,
wbrew dobrej redukcji. Dlatego $\bar\alpha_i$ są
$n$ różnymi pierwiastkami $\bar f$ w $\kappa$.
Ponieważ $\bar f$ ma stopień $n$, są to wszystkie pierwiastki
w dowolnym domknięciu algebraicznym $\mathbb F_p$.

Dla każdego $i$ mamy
$\overline{\sigma(\alpha_i)}=\bar\alpha_i^{\,p}$.
Tak więc bijekcja $\alpha_i\mapsto\bar\alpha_i$
przekształca permutację pierwiastków wywołaną przez $\sigma$
w permutację wywołaną przez Frobeniusa.
Na mocy Twierdzenia~\ref{aa:tw:frob-cycles} jej długości cykli
są dokładnie równe $d_1,\ldots,d_s$.
\end{proof}

\begin{uwaga}\label{aa:uw:dedekind-prereq}
W dowodzie użyto: twierdzenia Cayleya–Hamiltona dla macierzy
nad $\mathbb Z$, podstawowego twierdzenia teorii Galois,
cykliczności grupy multiplikatywnej ciała skończonego,
chińskiego twierdzenia o resztach oraz twierdzeń o pierścieniach
i ciałach ilorazowych. Kluczowe fakty o pierścieniu
$B=\mathbb Z[\alpha_1,\ldots,\alpha_n]$ i o podnoszeniu
Frobeniusa zostały powyżej udowodnione; nie jest wymagane
przyjmowanie skończoności całego pierścienia liczb całkowitych
$\mathcal O_L$ jako nieudowodnionego twierdzenia.
\end{uwaga}

\begin{wniosek}[Parzystość permutacji Frobeniusa]
\label{aa:wn:frob-parity}
Niech $f\in\mathbb Z[X]$ będzie monicznym, rozdzielczym
wielomianem stopnia $n\ge2$ i niech $p$ będzie liczbą pierwszą
niedzielącą $\operatorname{disc}(f)$. Załóżmy, że redukcja
$\bar f\in\mathbb F_p[X]$ ma dokładnie $s$ różnych monicznych
czynników nierozkładalnych, o stopniach $d_1,\ldots,d_s$.
Wtedy grupa Galois wielomianu $f$ nad $\mathbb Q$ zawiera
permutację pierwiastków o znaku $(-1)^{n-s}$. W szczególności,
jeśli $n-s$ jest nieparzyste, wyróżnik $\operatorname{disc}(f)$
nie jest kwadratem liczby wymiernej. Jeżeli wyróżnik jest kwadratem
liczby wymiernej, to dla każdej dobrej liczby pierwszej $p$
liczba $n-s$ musi być parzysta.
\end{wniosek}
\begin{proof}
Z Twierdzenia~\ref{aa:tw:dedekind-cycle} wynika, że istnieje
automorfizm $\sigma$ ciała rozkładu, który działa na zbiorze
pierwiastków jako iloczyn $s$ rozłącznych cykli długości
$d_1,\ldots,d_s$. Każdy cykl długości $d_i$ jest iloczynem
$d_i-1$ transpozycji, a więc ma znak $(-1)^{d_i-1}$.
Znak całej permutacji wynosi zatem
\[
 \operatorname{sgn}(\sigma)=\prod_{i=1}^s(-1)^{d_i-1}
 =(-1)^{d_1+\cdots+d_s-s}=(-1)^{n-s}.
\]
Jeśli $n-s$ jest nieparzyste, grupa Galois nie zawiera się
w $A_n$. Ponieważ $\operatorname{char}\mathbb Q=0$, z kryterium wyróżnika
(Twierdzenie~\ref{aa:tw:disc-alternating}) wynika, że grupa
zawiera się w $A_n$ dokładnie wtedy,
gdy wyróżnik jest kwadratem w $\mathbb Q$.
Stąd wyróżnik nie jest kwadratem. Ostatnie zdanie wynika
przez przeciwstawienie.
\end{proof}

\section{Przykłady i granice stosowalności}
\begin{przyklad}[Dwa rodzaje cykli]\label{aa:prz:modp-cycles}
Niech $f=X^3-X-1\in\mathbb Z[X]$.
Ze wzoru na wyróżnik wielomianu sześciennego
(Twierdzenie~\ref{aa:tw:cubic-disc}) mamy
$\operatorname{disc}(f)=-4(-1)^3-27(-1)^2=4-27=-23$.
Zatem zarówno $2$, jak i $5$ są liczbami dobrymi.

Modulo $2$ dostajemy $X^3+X+1$. Jego wartości dla
$0$ i $1$ wynoszą $1$, więc nie ma pierwiastków
w $\mathbb F_2$. Wielomian stopnia trzy bez pierwiastka
nad ciałem jest nierozkładalny, a więc otrzymujemy cykl
$(123)$ długości trzy.

Modulo $5$ mamy
\[
 \bar f=(X-2)(X^2+2X+3).
\]
Pierwszy czynnik jest liniowy, a drugi ma wyróżnik
$2^2-4\cdot3= -8\equiv2\pmod5$.
Kwadraty w $\mathbb F_5$ to $0,1,4$, więc czynnik
kwadratowy jest nierozkładalny. Otrzymujemy typ cykli
$(1,2)$, czyli transpozycję.
Dla $p=2$ mamy $n=3$, $s=1$, zatem znak otrzymanego
$3$-cyklu jest $(-1)^{3-1}=1$. Dla $p=5$ mamy $s=2$ i
znak permutacji o typie $(1,2)$ wynosi $(-1)^{3-2}=-1$.
Wniosek~\ref{aa:wn:frob-parity} potwierdza więc, że
wyróżnik $-23$ nie jest kwadratem liczby wymiernej.

Grupa Galois jest podgrupą $S_3$ zawierającą element
rzędu $3$ i element rzędu $2$, więc jej rząd
jest podzielny przez $6$. Jako dzielnik $6$ musi być
równy $6$. Zatem $\operatorname{Gal}(f/\mathbb Q)\cong S_3$.
\end{przyklad}

\begin{przyklad}[Dlaczego warunek dobrej redukcji jest konieczny]\label{aa:prz:bad-prime}
Dla $f=X^2-5$ wyróżnik wynosi $20$, więc liczba pierwsza
$2$ jest zła. Modulo $2$ mamy
$\bar f=X^2+1=(X+1)^2$.
Dwa pierwiastki nad $\mathbb Q$ redukują się więc do
jednego pierwiastka podwójnego. Nie wolno odczytywać
z zapisu $(X+1)^2$ cyklu długości dwa: wielomian
nie jest rozdzielczy po redukcji. Twierdzenie
\ref{aa:tw:dedekind-cycle} nie daje dla $p=2$ typu cyklu.

Tak samo dla wielomianu $g=X^3-X-1$ z
Przykładu~\ref{aa:prz:modp-cycles} liczba $23$ jest zła,
gdyż $23\mid\operatorname{disc}(g)=-23$.
W $\mathbb F_{23}[X]$ otrzymujemy
$\bar g=(X-3)(X-10)^2$. Rzeczywiście,
$(X-10)^2=X^2+3X+8$ modulo $23$, a pomnożenie
przez $X-3$ daje $X^3-X-1$.
Podwójny pierwiastek $10$ nie wyznacza cyklu długości dwa:
redukcja skleja tu dwa różne pierwiastki z charakterystyki zero.
\end{przyklad}

\begin{przyklad}[Rozkład na czynniki liniowe nie oznacza grupy trywialnej]
Niech $f=X^3-X-1$ jak wyżej. Weźmy $p=59$; liczba ta jest dobra, bo $59\nmid23$. Rachunek
modulo $59$ daje $\bar f(4)=\bar f(13)=\bar f(42)=0$.
Wszystkie trzy liczby są różne modulo $59$, więc
$\bar f=(X-4)(X-13)(X-42)$ w $\mathbb F_{59}[X]$.
Twierdzenie~\ref{aa:tw:dedekind-cycle}
dostarcza zatem w grupie $S_3$ elementu o typie $(1,1,1)$,
czyli elementu identycznościowego. Nie wynika z tego,
że cała grupa Galois jest trywialna: wiemy już,
że jest izomorficzna z $S_3$. Informacja z jednej liczby
pierwszej dotyczy istnienia \emph{pewnego} elementu,
a nie wszystkich elementów grupy.
\end{przyklad}

\begin{przyklad}[Wielomian rozkładalny nad $\mathbb Q$]
\label{aa:prz:frob-reducible}
Nie trzeba zakładać nierozkładalności $f$ nad $\mathbb Q$.
Rozważmy
$f=(X^2-2)(X-3)=X^3-3X^2-2X+6$.
Wyróżnik iloczynu dwóch monicznych wielomianów, których
zbiory pierwiastków są rozłączne, jest iloczynem
ich wyróżników i kwadratu wypadkowej. W naszym przypadku
$\operatorname{disc}(X^2-2)=8$, a
$R(X^2-2,X-3)=7$, więc
$\operatorname{disc}(f)=8\cdot 7^2=392$.
Dla $p=5$ wyróżnik nie znika modulo $5$, natomiast
$\bar f=(X^2-2)(X-3)$, gdzie $X^2-2$ jest
nierozkładalny nad $\mathbb F_5$, ponieważ $2$ nie
jest kwadratem w tym ciele. Otrzymujemy więc element
grupy Galois o typie cykli $(2,1)$.
Można to też zobaczyć bezpośrednio: ciało rozkładu
jest równe $\mathbb Q(\sqrt 2)$, a jego nietrywialny
automorfizm zamienia $\sqrt2$ z $-\sqrt2$ i ustala $3$.
\end{przyklad}
